\section{Soal}
Diberikan segitiga sama sisi $ABD$. Misalkan $C$ dan $O$ berturut-turut titik tengah $BD$ dan $AB$. Misalkan $I$ adalah pusat lingkaran dalam segitiga $ABC$ dan $N$ adalah pusat lingkaran luar segitiga $BIC$. Misalkan $\Gamma$ adalah lingkaran luar $ABC$. Lingkaran $\Omega$ menyinggung sisi $AC$, $BC$, dan menyinggung $\Gamma$ secara internal. Misalkan $M$ adalah titik pusat lingkaran $\Omega$. Misalkan $L$ adalah titik singgung lingkaran $\Omega$ dengan sisi $BC$ dan $X$ adalah titik singgung $\Omega$ dengan $\Gamma$. Buktikan bahwa garis $AB$, garis $CX$, lingkaran luar segitiga $BIC$, dan lingkaran luar segitiga $CNM$ berpotongan di satu titik.

\section{Motivasi dan Cerita}
    \begin{center}
       \begin{tikzpicture}[scale=1]
            % LANGKAH 1: Pengaturan Dasar
            % Karena O ada di AB, maka AB adalah diameter dan segitiga ABC siku-siku di C.
            \def\R{4} % Jari-jari lingkaran luar
            \tkzDefPoint(0,0){O}
            \tkzDefPoint(-\R,0){A}
            \tkzDefPoint(\R,0){B}
            \tkzDefPoint(60:\R){C}
        
            % Definisikan lingkaran luar Gamma
            \tkzDefCircle[R](O,\R) 
                \tkzGetPoint{gamma_pt}
            \tkzDrawCircle[thick](O,A)
            \tkzLabelCircle[above left=0.2 of A](O,gamma_pt)(180){$\Gamma$}
        
            % Gambar segitiga ABC
            \tkzDrawPolygon(A,B,C)
        
            % LANGKAH 2: Tentukan titik I dan N dan circumcircle BIC
            % I adalah pusat lingkaran dalam
            \tkzDefTriangleCenter[in](A,B,C)
                \tkzGetPoint{I}
            % N adalah pusat lingkaran luar segitiga BIC
            \tkzDefTriangleCenter[circum](B,I,C)
                \tkzGetPoint{N}
                \tkzLabelPoints[above right](N)
            % circumcircle BIC
            \tkzDrawCircle[teal](N,C)
            \tkzDrawSegment[teal](C,N)
            \tkzDrawSegment[teal, thick](I,N)
            \tkzDrawSegment[teal](B,N)
            \tkzDrawSegment[teal, thick](A,I)
            
            
        
            % LANGKAH 3: Konstruksi Titik M 
            % 3.1. Buat garis melalui I tegak lurus CI, lalu temukan D dan E = L
            \tkzDefLine[orthogonal=through I](C,I) \tkzGetPoint{p_perp}
            \tkzInterLL(I,p_perp)(A,C)
                \tkzGetPoint{D}
                \tkzDrawPoints(D)
            \tkzInterLL(I,p_perp)(B,C)
                \tkzGetPoint{L}
                \tkzDrawPoints(L)
            % \tkzDrawSegment[green](I,p_perp)
        
            % 3.2. Temukan titik X
            % Titik N (pusat lingkaran luar BIC) selalu terletak pada lingkaran Gamma.
            % Garis ND memotong Gamma di N dan titik kedua, X.
            \tkzInterLC(N,L)(O,A)
                \tkzGetPoints{X}{N1}
                \tkzDrawPoints[red](X)
                \tkzLabelPoints[red](X)
            \tkzDrawSegment[dashed](C,X)
            
            % 3.3. M adalah pusat lingkaran luar dari segitiga XDE (XDL)
            \tkzDefTriangleCenter[circum](X,D,L)
                \tkzGetPoint{M}
        
            % Gambar lingkaran Omega untuk verifikasi, seharusnya menyinggung AC, BC, dan Gamma
            \tkzDefPointBy[projection=onto A--C](M)
                \tkzGetPoint{P_ac}
            \tkzDrawCircle[color=blue](M,P_ac)
            \tkzLabelPoint[below=4 of D, color=blue](D){$\Omega$}
            
            % LANGKAH 4: Lingkaran luar MNC
            % Definisikan pusat lingkaran luar MNC
            \tkzDefTriangleCenter[circum](M,N,C)
                \tkzGetPoint{O_mnc}
            % Gambar lingkaran luar MNC (putus-putus)
            \tkzDrawCircle[red, dashed](O_mnc,C)
            
        
            % LANGKAH 5: Gambar Elemen Akhir dan Label
            % Gambar segmen 
            \tkzDrawSegments[red, thick, line cap=round](B,I O,L O,M)
            \tkzDrawSegment[thick](C,M)
        
            % Gambar titik-titik
            \tkzDrawPoints[size=4, fill=black](A,B,C,O,I,N,M,L)
        
            % Beri label pada semua titik
            \tkzLabelPoints[below left](A,B)
            \tkzLabelPoints[above right](C)
            \tkzLabelPoints[below, yshift=-3pt](O)
            \tkzLabelPoints[above right](I)
            
            \tkzLabelPoints[below right](M)
            \tkzLabelPoints[below](L)
        
        \end{tikzpicture} 
    \end{center}

    \textit{To be frank}, Ini adalah salah satu soal yang aku \textit{propose} ke Pak Barra untuk dijadikan kandidat soal OSN SMA 2025 tingkat Semifinal-Final. Namun, karena sepertinya terlalu sulit dan materinya melibatkan materi \textit{alien} berupa \textit{mixtilinear circles}, kayaknya ngga dijadiin soal OSN dari awal wkwkwkwk. Yaudahlah karena sayang udah dibuat susah payah, aku jadiin soal KTOM aja hehe.

    Anyway, aku kepikiran soal ini karena... baca bukunya Evan Chen dan kebayang beberapa properti yang bisa dimanipulasi pas ngerjain soal All Russian 2018 (kalo ngga salah). Intinya setelah main-main dengan Geogebra beberapa kali jadi deh soal ini. Ya intinya emang sih biar bisa nyadar beberapa properti harus paling ngga \textit{aware} kalo lingkaran \textit{mixti} itu \textit{exists} :).
    
    \newpage

    \section{Solusi}
    Jelas bahwa $\triangle ABC$ adalah segitiga siku-siku dengan $\angle ACB = 90^\circ$ dan $\angle ABC=60^\circ$. WLOG $AC > BC$. Perhatikan karena $O$ berada di $AB$ maka $AB$ adalah diameter lingkaran $\Gamma$ dengan $O$ adalah pusatnya.
    
    Misalkan $CM$ memotong $\Gamma$ untuk kedua kalinya di $T$. Definisikan $X$ sebagai titik singgung $\Gamma$ dan $\Omega$ serta $Y \neq X$ sebagai perpotongan $XM$ dengan $\Gamma$. Terakhir, misalkan $P$ dan $Q$ berturut-turut adalah titik singgung lingkaran dalam $\triangle ABC$ dengan sisi $BC$ dan $AC$.

    Misalkan panjang jari-jari $\Omega$ adalah $x$, panjang jari-jari lingkaran luar $\triangle BIC$ adalah $R$, panjang jari-jari lingkaran dalam $\triangle ABC$ adalah $r$ dan panjang sisi $AB=c$, $BC=a$, $CA=b$.

    \begin{center}
        \begin{tikzpicture}[scale=1]
            % LANGKAH 1: Pengaturan Dasar
            % Karena O ada di AB, maka AB adalah diameter dan segitiga ABC siku-siku di C.
            \def\R{4} % Jari-jari lingkaran luar
            \tkzDefPoint(0,0){O}
            \tkzDefPoint(-\R,0){A}
            \tkzDefPoint(\R,0){B}
            \tkzDefPoint(60:\R){C} %% B 60 degree
        
            % Definisikan lingkaran luar Gamma
            \tkzDefCircle[R](O,\R) 
                \tkzGetPoint{gamma_pt}
            \tkzDrawCircle[thick](O,A)
            \tkzLabelCircle[above left=0.2 of A](O,gamma_pt)(180){$\Gamma$}
        
            % Gambar segitiga ABC
            \tkzDrawPolygon(A,B,C)
        
            % LANGKAH 2: Tentukan titik I dan N dan circumcircle BIC
            % I adalah pusat lingkaran dalam
            \tkzDefTriangleCenter[in](A,B,C)
                \tkzGetPoint{I}
            % N adalah pusat lingkaran luar segitiga BIC
            \tkzDefTriangleCenter[circum](B,I,C)
                \tkzGetPoint{N}
                \tkzLabelPoints[above right](N)
            % circumcircle BIC
            \tkzDrawCircle[teal](N,C)
            \tkzDrawSegment[teal, thick](C,N)
            \tkzDrawSegment[teal, thick](I,N)
            \tkzDrawSegment[teal, thick](B,N)
            \tkzDrawSegment[teal, thick](A,I)
            
            % LANGKAH 3: Konstruksi Titik M 
            % 3.1. Buat garis melalui I tegak lurus CI, lalu temukan D dan E = L
            \tkzDefLine[orthogonal=through I](C,I) \tkzGetPoint{p_perp}
            \tkzInterLL(I,p_perp)(A,C)
                \tkzGetPoint{D}
            \tkzInterLL(I,p_perp)(B,C)
                \tkzGetPoint{L}
                \tkzDrawPoints(L)
            % \tkzDrawSegment[green](I,p_perp)
        
            % 3.2. Temukan titik X
            % Titik N (pusat lingkaran luar BIC) selalu terletak pada lingkaran Gamma.
            % Garis ND memotong Gamma di N dan titik kedua, X.
            \tkzInterLC(N,L)(O,A)
                \tkzGetPoints{X}{N1}
                \tkzDrawPoints[red](X)
                \tkzLabelPoints[red](X)
            
            % 3.3. M adalah pusat lingkaran luar dari segitiga XDE (XDL)
            \tkzDefTriangleCenter[circum](X,D,L)
                \tkzGetPoint{M}
        
            % Gambar lingkaran Omega untuk verifikasi, seharusnya menyinggung AC, BC, dan Gamma
            \tkzDrawCircle[color=blue](M,D)
            \tkzLabelPoint[below=4 of D, color=blue](D){$\Omega$}
            
            % LANGKAH 4: Lingkaran luar MNC
            % Definisikan pusat lingkaran luar MNC
            \tkzDefTriangleCenter[circum](M,N,C)
                \tkzGetPoint{O_mnc}
            % Gambar lingkaran luar MNC (putus-putus)
            \tkzDrawCircle[red, dashed](O_mnc,C)

            \tkzDrawSegment[teal](M,N)
            
        
            % LANGKAH 5: Gambar Elemen Akhir dan Label
            % Gambar segmen 
            \tkzDrawSegment[red, line cap=round](B,I)
            \tkzDrawSegment[red, thick, line cap=round](M,L)
            \tkzDrawSegment[thick](C,M)

            % CT perpotongan CM dengan Gamma
            \tkzInterLC(C,M)(O,A)
                \tkzGetPoints{T}{C}
            \tkzLabelPoints[below right](T)
            \tkzDrawSegment[thick](C,T)

            % proyeksi I ke BC dan AC
            \tkzDefPointBy[projection=onto B--C](I)
                \tkzGetPoint{P}
            \tkzDrawPoint[fill=black](P)
            \tkzLabelPoints[right](P)
            \tkzDrawSegment[thick, red](I,P)

             \tkzDefPointBy[projection=onto A--C](I)
                \tkzGetPoint{Q}
            \tkzDrawPoint[fill=black](Q)
            \tkzLabelPoints[left](Q)
            \tkzDrawSegment[red](I,Q)

            % AT = BT
            \tkzDrawSegments[orange, thick](A,T B,T)

            % XMO cuts Gamma at Y
            \tkzInterLC(X,M)(O,A)
                \tkzGetPoints{X}{Y}
                \tkzDrawPoints[red](Y)
                \tkzLabelPoints[above left, red](Y)
            \tkzDrawSegment(X,Y)

            % CX, OC, LK
            \tkzInterLL(C,X)(B,A)
                \tkzGetPoint{K}
                \tkzDrawPoints(K)
                \tkzLabelPoints[below](K)
            \tkzDrawSegments(C,X L,K O,L O,C M,K)
            \tkzDrawSegments[red, thick](O,K C,L C,D)

            % N adalah pusat lingkaran luar segitiga BIC
            \tkzDefTriangleCenter[circum](B,I,C)
                \tkzGetPoint{N}
                \tkzLabelPoints[above right](N)
            % circumcircle BIC
            \tkzDrawCircle[teal](N,C)

            \tkzDrawSegments[](N,O O,D D,M)

            % perpotongan CX OL
            \tkzInterLL(O,L)(C,X)
                \tkzGetPoint{J}
                \tkzLabelPoints[above left](J)
                \tkzDrawPoints(J)
                
            % Gambar titik-titik
            \tkzDrawPoints[size=4, fill=black](A,B,C,O,I,N,M,L,D,T)
        
            % Beri label pada semua titik
            \tkzLabelPoints[below left](A,B)
            \tkzLabelPoints[above right](C)
            \tkzLabelPoints[below, yshift=-3pt](O)
            \tkzLabelPoints[above right](I)
            
            \tkzLabelPoints[below right](M)
            \tkzLabelPoints[below](L)
            \tkzLabelPoints[above left](D)
        
        \end{tikzpicture}
    \end{center}

    \newpage
    \begin{proof}[\textbf{Klaim 1.}] Panjang $r = \dfrac{a+b-c}{2}$.\\
        \textbf{Bukti Klaim.} Perhatikan bahwa $\angle C = 90^\circ$ sehingga $r = IP = CP = s-c = \dfrac{a+b-c}{2}$ dimana $s = \dfrac{a+b+c}{2}$. Klaim 1 terbukti.
    \end{proof}
    \hrulefill
    \begin{proof}[\textbf{Klaim 2.}] Panjang $x = a+b-c$.\\
        \textbf{Bukti Klaim.} Perhatikan bahwa terdapat sebuah homothety atau dilatasi berpusat di $X$ yang mengirim lingkaran $\Omega$ ke lingkaran $\Gamma$ karena kedua lingkaran tersebut bersinggungan di $X$. Dari sini, jelas bahwa $M$ dikirim ke $O$ sehingga $X,M,O,Y$ segaris yang mana menunjukkan juga bahwa $XY$ diameter $\Gamma$. 
        
        Lalu, karena $AB$ diameter $\Gamma$, maka $\angle ACB = 90^\circ$ sehingga $\angle MCL = \frac{1}{2}\angle ACB = 45^\circ$. Karena $ML = x$ jari-jari $\Omega$, maka $CM = x\sqrt{2}$. Jangan lupa fakta kalau $MX = x$ dan $XY=AB=c$. Oleh karena itu dengan \textit{power of a point} didapat
        \begin{align*}
            CM \cdot MT &= XM \cdot MY\\
            (CT - CM) \cdot CM &= XM \cdot (XY - XM)\\
            (\frac{\sqrt{2}}{2}(a+b) - x\sqrt{2}) \cdot x\sqrt{2} &= x \cdot (c - x)\\
            (a+b)x - 2x^2 &= cx - x^2\\
            a+b - 2x &= c - x\\
            x &= a+b -c
        \end{align*}
        Klaim 2 Terbukti.
    \end{proof}
    
    \newpage
    
    Dari klaim 1 dan klaim 2 kita punya $\dfrac{IP}{ML} = \dfrac{r}{x} = \dfrac{1}{2}$. Sekarang, karena $IP \perp BC$ dan $ML \perp BC$ maka $IP \parallel ML$ sehingga $\triangle CIP \sim \triangle CML$ yang menyebabkan
    \begin{align*}
        \dfrac{CI}{CM} = \dfrac{IP}{ML} = \dfrac{1}{2} \implies CI = IM
    \end{align*}

    \hrulefill
    \begin{proof}[\textbf{Lemma 1. (Incenter-Excenter Lemma)}] (Tidak wajib dibuktikan untuk saat kontes berlangsung) $A,I,N$ segaris dengan $N$ berada di lingkaran $\Gamma$\\
        \textbf{Bukti Lemma.} Misalkan $N_1 \neq A$ adalah perpotongan $AI$ dengan $\Gamma$. Akan dibuktikan bahwa $N_1 = N$. 
        Perhatikan karena $AI$ garis bagi $\angle CAB$, maka $$\angle N_1BC = \angle N_1AC = \angle BAN_1 = \angle BCN_1 \implies N_1C=N_1B.$$ 
        Lalu observasi bahwa
        \begin{align*}
            \angle N_1IC &= 180^\circ - \angle CIA = \angle ACI + \angle IAC = \frac{1}{2}\angle ACB + \frac{1}{2}\angle BAC\\ &= 90^\circ - \frac{1}{2}\angle CBA = 90^\circ - \frac{1}{2} CN_1A
        \end{align*}
        sehingga didapat
        $$\angle ICN_1 = 180^\circ - \angle N_1IC - \angle CN_1A = 180^\circ - \angle N_1IC - (180^\circ - 2\angle N_1IC) = \angle N_1IC$$
        yang menyebabkan $N_1I=CN_1$. Dari sini didapat $N_1B=N_1C=N_1I$, yang menunjukkan bahwa $N_1$ adalah pusat lingkaran luar $\triangle BIC$ sehingga $N_1 = N$. Oleh karena itu, $A,I,N$ segaris dengan $N$ berada di lingkaran $\Gamma$. Lemma terbukti.
    \end{proof}
    \newpage

    \begin{proof}[\textbf{Klaim 3.}] $AD=AO$.\\
        \textbf{Bukti.} Karena $\angle ABC = 60^\circ$ dan $\angle ACB = 90^\circ$ maka $AB = 2\cdot BC \implies c = 2a$.

        Perhatikan bahwa $CD = CL = ML = DM = x$. Dari Klaim 2 kita punya
        \begin{align*}
            AD = AC - CD = b - x = b - (a+b-c) = c - a = a = \frac{1}{2}AB = AO
        \end{align*}
        Klaim terbukti.
    \end{proof}
    \hrulefill

    \begin{proof}[\textbf{Klaim 4.}] $M,D,C,N,L$ berada di satu lingkaran yang sama.\\
        \textbf{Bukti.} Karena $\angle ABC = 60^\circ$, dari Lemma 1, dan fakta bahwa $CI=IM$, maka $\angle CNI = \angle CNA = \angle CBA = 60^\circ \text{ dan } NC = NI$ sehingga didapat  $\triangle CNI$ sama sisi yang menyebabkan $MI = IC = IN$. Oleh karena itu $MC = 2 \cdot CN$  dan $\angle MCN = 60^\circ$ sehingga didapat $\angle MNC = 90^\circ$. Dari definisi $D$ dan $L$, kita juga punya $\angle MDC = \angle MLC = 90^\circ$. Dari sini didapat $M,D,C,N,L$ berada di satu lingkaran yang sama. Terbukti.
    \end{proof}
    \hrulefill

    \begin{proof}[\textbf{Klaim 5.}] $C,D,O,N$ berada di satu lingkaran yang sama.\\
        \textbf{Bukti.} Karena $CO=ON$ dan $N$ titik tengah busur $BC$ maka $\angle NOC = \frac{1}{2}\angle BOC = \frac{1}{2}\angle 60^\circ = 30^\circ$. Dari sini didapat juga $\angle CNO = 75^\circ$. 
        
        Lalu, karena $AD=AO$ dari klaim 3 dan $\angle OAD = 30^\circ$ maka $\angle ODA = 75^\circ \implies \angle CDO = 180^\circ - 75^\circ = 180^\circ - \angle CNO$. Diperoleh $C,D,O,N$ berada di satu lingkaran yang sama. Terbukti
    \end{proof}
    \hrulefill

    Dari fakta-fakta klaim 4 dan klaim 5 tersebut diperoleh $M,D,C,L,N,O$ berada di satu lingkaran yang sama.


    Karena $\angle COX = 90^\circ$ dan $CO=CX$ (jari-jari $\Gamma$), maka $\angle OCX = 45^\circ$. Berarti $\angle DCX = \angle ACO + \angle OCX = 30^\circ + 45^\circ = 75^\circ$.

    Misalkan lingkaran $(DCLMO)$ memotong $AB$ sekali lagi di $K$. Berarti $\angle CKO = 180^\circ - \angle CDO = \angle ADO = 75^\circ$. Lalu, $\angle COK = \angle COB = 60^\circ$. Dari sini, didapat $\angle OCK = 180-\angle COK - \angle CKO = 45^\circ = \angle OCX$. Diperoleh $C,K,X$ segaris yang menyebabkan lingkaran luar $\triangle CNM$, $AB$, $CX$ bertemu di satu titik.

    Namun, kita juga tahu bahwa $\angle CNB = 180^\circ - \angle BAC = 150^\circ$. Berarti sudut refleks $\angle CNB = 360^\circ - 150^\circ = 210 ^\circ = 2\cdot 105^\circ = 2(180^\circ - 75^\circ) = 2(180^\circ - \angle CKO) = 2\angle AKB$. Terbukti bahwa lingkaran luar $\triangle BIC$ yang berpusat di $N$ melewati $K$.

    Terbukti bahwa garis $AB$, garis $CX$, lingkaran luar $\triangle BIC$, dan lingkaran luar $\triangle CNM$ bertemu di satu titik. 

    \newpage

    \section{Marking Scheme}
	\begin{enumerate}[1)]
	\item Mendapatkan fakta $M,D,C,L,O$ berada di satu lingkaran yang sama. (+5 poin)
    \begin{enumerate}[(a)]
        \item Memperoleh $CI=IM$ (+2 poin)
        \begin{enumerate}[(i)]
            \item Memperoleh Klaim 1 dan Klaim 2 (memperoleh panjang jari-jari lingkaran dalam $ABC$ dan lingkaran $\Omega$) (+1 poin)
            \item Melengkapkan solusi sehingga $CI=IM$ (+1 poin)
        \end{enumerate}
        \item Melengkapi solusi agar didapat $M,D,C,L,O$ konsiklis (+3 poin)
        \begin{enumerate}[(i)]
            \item Memperoleh Klaim 3 dan memakainya di Klaim 5 (+1 poin)
            \item Memperoleh Klaim 4 dengan menyebut Lemma 1 (Lemma 1 tidak wajib dibuktikan, dianggap well known), serta memakai $CI=IM$ (+2 poin)
        \end{enumerate}
    \end{enumerate}
	\item Membuktikan atau melengkapkan solusi: lingkaran luar $\triangle BIC$ yang berpusat di $N$ melewati $K$. (+2 poin)
    \begin{enumerate}[(a)]
        \item Memperoleh $C,K,X$ segaris yang menyebabkan lingkaran luar $\triangle CNM$, $AB$, $CX$ bertemu di satu titik. (+1 poin)
        \item Memperoleh lingkaran luar $\triangle BIC$ yang berpusat di $N$ melewati $K$. (+1 poin)
    \end{enumerate}
	\end{enumerate}